% ============================================================
%  part3/ch12-ising.tex  —  Chapter 12: The Ising Model
% ============================================================

\chapter{The Ising Model}
\label{ch:ising}

\epigraph{%
  The simple model already contains the germ of everything.%
}{---}

\begin{WhyMatters}
  The Ising model is introduced not because the universe is
  literally an Ising lattice, but because it is the simplest
  system exhibiting: local interactions, domain formation,
  symmetry breaking, coarsening, and defect stabilisation.
  These are precisely the ingredients later required for
  RSVP cosmology.
  Every phenomenon in Chapters~13--18 and every structure in
  Part~VII is a generalisation of what first appears here.
\end{WhyMatters}

\section{Setup}

\begin{definition}[Ising model]
  Let $\Lambda \subset \mathbb{Z}^d$ be a lattice.
  Each site $i \in \Lambda$ carries a spin $\sigma_i \in \{-1,+1\}$.
  The Hamiltonian is
  \[
    H(\sigma) = -J \sum_{\langle i,j\rangle} \sigma_i\sigma_j,
    \quad J > 0,
  \]
  where the sum runs over nearest-neighbour pairs.
\end{definition}

\begin{proposition}
  Aligned states minimise energy.
\end{proposition}

\begin{proof}
  An aligned pair contributes $-J$; an anti-aligned pair $+J$.
  Minimum energy is achieved when all pairs are aligned.
\end{proof}

\begin{theorem}[Ground states]
  \label{thm:ising-ground}
  The Ising Hamiltonian possesses exactly two global minima:
  $\sigma_i \equiv +1$ and $\sigma_i \equiv -1$.
\end{theorem}

\begin{proof}
  Every neighbouring pair must align (Proposition above).
  Connectivity of $\Lambda$ forces all spins to match.
  Since $\sigma_i \in \{-1,+1\}$, the two uniform states are
  the only ground states.
\end{proof}

This is the first appearance of \emph{symmetry breaking}: the
Hamiltonian is symmetric under $\sigma \mapsto -\sigma$, but
each ground state breaks this symmetry.

\section{Domain walls}

When regions of opposite sign coexist, a \emph{domain wall}
separates them.

\begin{proposition}[Domain wall energy]
  \label{prop:wall-energy}
  Domain wall energy scales as surface area:
  $E_{\mathrm{wall}} \propto |\partial D|$.
\end{proposition}

\begin{proof}
  Each broken bond contributes $2J$ above the ground state.
  The number of broken bonds equals the number of
  nearest-neighbour pairs straddling the wall, which scales
  with $|\partial D|$.
\end{proof}

\principle{Interfaces are energetically expensive.\\[4pt]
  The rest of Part~III is the mathematics of removing them.}

\section{Curvature-driven motion}

A bulging interface has positive curvature $\kappa = 1/R$.
Reducing curvature lowers interface energy.

\textbf{Heuristic derivation.}
Interface energy: $E = \gamma|\partial D|$.
First variation: $\delta E = -\gamma \int_{\partial D} \kappa V_n\, dS$.
Gradient descent requires:
\[
  \boxed{V_n \propto -\kappa.}
\]
Small structures shrink faster than large ones.

\begin{theorem}[Characteristic scale growth]
  \label{thm:ising-scaling}
  Under curvature-driven motion, the characteristic domain
  size grows as $\ell(t) \sim t^{1/2}$.
\end{theorem}

\begin{proof}
  Since $\kappa \sim 1/\ell$, the interface velocity
  $V_n \sim 1/\ell$.
  Then $d\ell/dt \sim 1/\ell$, giving $\ell\, d\ell \sim dt$,
  so $\ell^2 \sim t$, hence $\ell(t) \sim t^{1/2}$.
\end{proof}

This is the first appearance of the coarsening exponent $z = 2$.

\begin{CrossDomain}
  \begin{center}
  \begin{tabular}{@{}lll@{}}
    \toprule
    \textbf{System} & \textbf{Order parameter} & \textbf{Exponent $z$} \\
    \midrule
    Ising magnet & Spin $\sigma_i$ & 2 \\
    Allen--Cahn & Scalar field $\phi$ & 2 \\
    Cahn--Hilliard & Conserved density & 3 \\
    RSVP (conjectured) & $(\Phi,\mathbf{v},S)$ & $z(Pe,\Gamma,\Omega)$ \\
    \bottomrule
  \end{tabular}
  \end{center}
\end{CrossDomain}

\WhatSurvived{Ising coarsening $\sigma(0) \mapsto \sigma(t)$}{%
  Domain count (decreasing)\\
  Ground-state symmetry\\
  Topological defects%
}{%
  Initial spin configuration\\
  Short-wavelength fluctuations\\
  Sub-domain-scale detail%
}{%
  Partial; long-time state\\
  is one of two ground states,\\
  but path is irreversible.%
}{%
  $\ker$: all initial conditions\\
  flowing to the same\\
  final domain state.%
}

\begin{ProblemSet}
\problem{
  Verify that the Ising Hamiltonian $H(\sigma) = -J\sum_{\langle i,j\rangle}\sigma_i\sigma_j$
  is invariant under the global spin flip $\sigma_i \mapsto -\sigma_i$.
  Show that the two ground states $\sigma_i \equiv +1$ and
  $\sigma_i \equiv -1$ are related by this symmetry.
}
\problem{
  For a 1D Ising chain of $N$ sites with periodic boundary conditions,
  count the number of domain walls as a function of the number of
  sign changes in $(\sigma_1, \ldots, \sigma_N)$.
  What is the minimum energy above the ground state that contains
  exactly two domain walls?
}
\problem{
  Using the scaling argument $d\ell/dt \sim 1/\ell$,
  show that $\ell(t) \sim t^{1/2}$.
  What initial condition $\ell(0)$ would give $\ell(1) = 1$ Mpc
  if time is measured in Gyr?
}
\problem{
  \textbf{Reconstruction exercise.}
  Suppose you observe the state of a coarsened Ising system
  at time $t = 100$.
  Domain size is $\ell = 10$.
  Can you determine the initial condition?
  What information about the initial state survives in the
  observed domain structure?
  What is lost?
  (This is a concrete instance of $E(x) = d(x, G(F(x))) > 0$.)
}
\end{ProblemSet}
